% Part: incompleteness
% Chapter: introduction
% Section: overview

\documentclass[../../../include/open-logic-section]{subfiles}

\begin{document}

\olfileid{inc}{int}{ovr}

\olsection{અપૂર્ણતાનાં પરિણામોનું અવલોકન}

હિલ્બર્ટને આશા હતી કે ગણિતને એવા
!!{axiomatizable} સિદ્ધાંતમાં ઔપચારિક બનાવી
શકાય જેને !!{complete} અને !!{decidable}
છે એમ સાબિત કરી શકાય. વધુમાં, તે આ
સિદ્ધાંતની સુસંગતતા ખૂબ નબળી, ``સાન્ત''
પદ્ધતિઓથી સાબિત કરવા ઇચ્છતા હતા;
તેના દ્વારા શાસ્ત્રીય ગણિતને સ્ફુરણાવાદના
પડકારોથી બચાવી શકાય. ગડેલનાં અપૂર્ણતાનાં
પ્રમેયોએ બતાવ્યું કે આ હેતુઓ હાંસલ થઈ
શકતા નથી.

ગડેલના પ્રથમ અપૂર્ણતા-પ્રમેયે બતાવ્યું કે
રસેલ અને વ્હાઇટહેડના \emph{Principia
Mathematica}નું એક સ્વરૂપ !!{complete}
નથી. પરંતુ સાબિતી ઘણી સામાન્ય હતી અને
વિવિધ પ્રકારના સિદ્ધાંતોને લાગુ પડે છે. આનો
અર્થ કે \emph{Principia Mathematica}
ગણિતને સંપૂર્ણ રીતે સમાવી શક્યું નહીં એટલું
જ નહીં, પરંતુ \emph{કોઈ પણ} સ્વીકાર્ય
સિદ્ધાંત તેમ કરી શકતો નથી. અપૂર્ણતાનાં
પ્રમેયો લાગુ પડવા માટે સિદ્ધાંતનાં કયાં
લક્ષણો પૂરતાં છે તે અલગ તારવવામાં, અને
ગડેલની સાબિતીને માત્ર આ લક્ષણો પર આધારિત
કરી સામાન્ય બનાવવામાં સમય લાગ્યો.
હવે આપણે અંકગણિતની ભાષા~$\Lang L_A$માં
રચાયેલા સિદ્ધાંતો માટે પ્રથમ અપૂર્ણતા-પ્રમેયનું
ખૂબ સામાન્ય સ્વરૂપ કહી શકીએ છીએ.

\begin{thm}
જો $\Gamma$ એ~$\Lang L_A$માં સુસંગત
અને !!{axiomatizable} સિદ્ધાંત હોય, જે બધાં
સંગણનીય વિધેયો અને !!{decidable} સંબંધોને
!!{represents} કરતો હોય, તો $\Gamma$
!!{complete} નથી.
\end{thm}

$\Gamma$ !!{complete} નથી એનો અર્થ એ
કે ઓછામાં ઓછું એક !!{sentence}~$!A$ એવું
છે કે $\Gamma \Proves/ !A$ અને
$\Gamma \Proves/ \lnot !A$. આવું
!!a{sentence} ($\Gamma)$થી \emph{સ્વતંત્ર}
કહેવાય છે. ખરેખર, સ્વતંત્ર વાક્યો અસ્તિત્વ
ધરાવે છે તે આપણે તદ્દન ઝડપથી સાબિત કરી
શકીએ. પરંતુ ગડેલની આ પ્રમેયની સાબિતીની
શક્તિ એ છે કે તે આવા સ્વતંત્ર !!{sentence}નું
એક \emph{ચોક્કસ ઉદાહરણ} આપે છે. આ
રસપ્રદ રચનાથી~$\Gamma$ માટેનું
!!a{sentence}~$!G_\Gamma$ મળે છે, જેને
\emph{ગડેલ-વાક્ય} કહે છે. તે સાબિત ન થઈ
શકે, કારણ કે~$\Gamma$માં $!G_\Gamma$
એવું કહેતા વિધાનને સમકક્ષ છે કે
$!G_\Gamma$ એ~$\Gamma$માં સાબિત
થઈ શકતું નથી. આ રચના \emph{રચનાત્મક}
રીતે થાય છે: $\Gamma$નું સ્વયંસિદ્ધીકરણ
અને !!{derivation} તંત્રનું વર્ણન આપવામાં
આવે તો સાબિતી $!G_\Gamma$ ખરેખર લખવાની
પદ્ધતિ આપે છે.

ગડેલની સાબિતીની રચનામાં~$\Lang L_A$ના
પદો અને !!{formula}sના ગુણધર્મો તથા તેમના
પરની ક્રિયાઓને $\Lang L_A$માં જ વ્યક્ત કરવાની
રીત જોઈએ. તેમાં ``$!A$ એ !!a{sentence}
છે'', ``$\delta$ એ~$!A$ની
!!a{derivation} છે'' જેવા ગુણધર્મો અને
$\Subst{!A}{t}{x}$ જેવી ક્રિયાઓ આવે છે.
આ રીતમાં (ક) ચિહ્નો અને તેમના અનુક્રમો—જેમાંથી
પદો, !!{formula}s, !!{derivation}s વગેરે
બને છે—તેમને પ્રાકૃતિક સંખ્યાઓમાં ``સંકેતબદ્ધ''
કરી આ ગુણધર્મો અને સંબંધો વ્યક્ત કરવા પડે;
કારણ કે~$\Lang L_A$ પ્રાકૃતિક સંખ્યાઓ
વિશે જ વાત કરી શકે છે. (ખ) સંકેતબદ્ધીકરણ
એવું હોવું જોઈએ કે~$\Gamma$ સંબંધિત તથ્યો
સાબિત કરે. તેથી બતાવવું પડે કે આ ગુણધર્મો
પ્રાકૃતિક સંખ્યાઓના !!{decidable} ગુણધર્મો
વડે અને ક્રિયાઓ પ્રાકૃતિક સંખ્યાઓ પરના
સંગણનીય વિધેયો વડે સંકેતબદ્ધ થાય છે.
આને ``વાક્યરચનાનું અંકગણિતીકરણ'' કહે છે.

વાક્યરચનાનું અંકગણિતીકરણ કેવી રીતે થાય છે
તે તપાસતાં પહેલાં, $\Gamma$ ``પૂરતો
શક્તિશાળી'' છે—એટલે બધાં સંગણનીય વિધેયો
અને !!{decidable} સંબંધોને !!{represents}
કરે છે—એ શરત જોઈએ. તે માટે
``સંગણનીય''ની ચોક્કસ વ્યાખ્યા આપવી પડશે.
આ અનેક રીતે કરી શકાય; દાખલા તરીકે,
ટ્યુરિંગ મશીનોના નિદર્શથી અથવા કોઈ
સામાન્ય ઉપયોગની પ્રોગ્રામિંગ ભાષામાંના
પ્રોગ્રામો વડે ગણાતાં વિધેયો તરીકે. પરંતુ
આ વિધેયો અને સંબંધોને ભાષા~$\Lang L_A$માં
રચાયેલા સિદ્ધાંતમાં !!{represents}s કરવાનો
આપણો હેતુ હોવાથી, માત્ર સંખ્યાત્મક વિધેયોની
સંગણનીયતાની સરળ વ્યાખ્યા પસંદ કરવી સારી.
તે \emph{પુનરાવર્તી વિધેય}નો ખ્યાલ છે.
તેથી પહેલાં આપણે પુનરાવર્તી વિધેયોની ચર્ચા
કરીશું. ત્યાર પછી બતાવીશું કે~$\Th{Q}$
પહેલેથી જ બધાં પુનરાવર્તી વિધેયો અને
સંબંધોને !!{represents} કરે છે. આમ
$\Th{Q}$ અને~$\Th{PA}$ જેવા ચોક્કસ
સિદ્ધાંતોને અપૂર્ણતા-પ્રમેય લાગુ કરી શકીશું,
કારણ કે તેઓ ``પૂરતા શક્તિશાળી'' સિદ્ધાંતોનાં
ઉદાહરણો છે તે સ્થાપિત થઈ જશે.

વાક્યરચનાના અંકગણિતીકરણનું અંતિમ પરિણામ
!!a{formula} $\OProv[\Gamma](x)$ છે. તે
!!{formula}sને સંખ્યાઓથી સંકેતબદ્ધ કરીને
$\Gamma$નાં સ્વયંસિદ્ધ વાક્યોમાંથી સાબિત
થવાની શક્યતા વ્યક્ત કરે છે. ચોક્કસ રીતે,
$!A$નો સંકેત નંબર~$n$ હોય અને
$\Gamma \Proves !A$ હોય, તો
$\Gamma \Proves \Prov[\Gamma](\num{n})$.
$\Gamma$ માટેનું આ ``સાબિતીક્ષમતા-વિધેય''
આપણને ચોક્કસ અર્થમાં~$\Gamma$ની સુસંગતતા
$\Lang L_A$ના !!a{sentence} તરીકે વ્યક્ત
કરવા દે છે. $\Gamma$ માટેનું
``સુસંગતતા-વિધાન'' !!{sentence}
$\lnot \Prov[\Gamma](\num{n})$ ગણીએ,
જ્યાં~$n$ કોઈ વિરોધાભાસનો સંકેત નંબર,
દાખલા તરીકે~$\lfalse$નો, છે. બીજો
અપૂર્ણતા-પ્રમેય કહે છે કે યોગ્ય વધુ શરતો
સંતોષતા સુસંગત !!{axiomatizable} સિદ્ધાંતો પોતાના
સુસંગતતા-વિધાનને પણ સાબિત કરતા નથી.
આ પ્રમેય લાગુ કરવા માટે માત્ર બધાં
સંગણનીય વિધેયો અને !!{decidable}
સંબંધોનું નિરૂપણ કરવાની શરત કરતાં
થોડી વધુ કડક શરતો જોઈએ; $\Th{PA}$
તે સંતોષે છે તે આપણે બતાવીશું.
\sourcecorrection{OLINC-003}{સ્થિર અંગ્રેજી મૂળના
પહેલા વાક્યમાં બધા સુસંગત સ્વયંસિદ્ધીકરણીય
સિદ્ધાંતો માટે પોતાની સુસંગતતા સાબિત ન થવાનો
દાવો આવે છે. એ જ ફકરો તરત વધુ કડક શરતો
જરૂરી હોવાનું કહે છે. તેથી અહીં દાવામાં
એ શરતો સ્પષ્ટ કરી છે.}

\end{document}
